% kernel/representations-DE.tex

\chapter{Den Kandidaten repräsentieren, ohne ihn zu ersetzen}
\label{chap:representations}

Der vorgeschlagene Übergang wird nun als stetige Kurve für die Berechnung, als
abgetastete Polylinie für die CAD-Prüfung und als Austauschdatensatz für die
Koordination erzeugt. Diese Formen machen den Kandidaten nutzbar; keine wird
zum Alignment oder besitzt seine Identität.

\begin{definition}[Repräsentation]
\label{def:representation}
Eine Repräsentation ist eine abgeleitete Beschreibung von
Ingenieurinformationen, die zur Kommunikation, Visualisierung, zum Austausch,
zur Berechnung oder Interaktion bestimmt ist.
\end{definition}

Jede Repräsentation muss den ausgedrückten Zustand, ihren Zweck, die erzeugende
Ableitung und jede Näherung angeben. Die CAD-Polylinie kann für visuelle
Koordination genügen und zugleich die Funktionsfamilienidentität und
Ableitungen auslassen, die zur Rekonstruktion der Berlinish-Komposition
erforderlich sind. Das stetige Berechnungsmodell kann die Krümmungsbewertung
tragen und als Austauschvertrag ungeeignet sein. Eignung ist zweckrelativ.

Transformation, Neuabtastung oder Neuerzeugung können deshalb eine weitere
Repräsentation desselben Kandidatenzustands hervorbringen. Umgekehrt müssen
zwei geometrisch gleiche Dateien nicht dasselbe Objekt betreffen.
Repräsentationsgleichheit ist weder Objektidentität noch Ingenieurannahme.

Ein importierter historischer Übergang verdeutlicht die Gegenrichtung. Seine
Datei kann innerhalb ihres Geltungsbereichs autoritative Quellevidenz bleiben,
auch wenn ihre abgetastete Geometrie nicht mathematisch optimal ist. Die
Ableitung eines editierbaren stetigen Modells bewahrt den Import und
protokolliert eine neue Ableitung; sie macht das abgeleitete Modell nicht
rückwirkend zur ursprünglichen Beobachtung.

Der Kandidat ist nun kommunizierbar. Für den Vergleich mit der vorgefundenen
Eisenbahn muss die Beschreibung als Nächstes erzeugte Geometrie von
beobachteter Evidenz unterscheiden.
